\documentclass[10pt,a4paper]{amsart}
\usepackage[margin=19mm]{geometry}
\usepackage{amsmath,amssymb,mathtools}
\usepackage[scr=rsfs,cal=euler]{mathalfa}
\usepackage{xfrac}
\usepackage[unicode,hidelinks,bookmarks=false]{hyperref}
\newcommand{\ie}{i.e.\ }
\newcommand{\cf}{cf.\ }
\newcommand{\resp}{respectively\ }
\newcommand{\Subsection}{\subsection}
\providecommand{\nobreakdash}{\nobreak\hbox{-}}
\setlength{\emergencystretch}{2em}
\multlinegap0pt
\usepackage[shortlabels]{enumitem}
\usepackage{mathtools}
\usepackage{bm}
\usepackage{csquotes}
\usepackage{amssymb}
\usepackage{tikz}
\newtheorem{definition}{Definition}
\newtheorem{theorem}{Theorem}
\newtheorem{corollary}{Corollary}
\newcommand{\A}{\boldsymbol A}
\newcommand{\Nzero}{\mathbb N_0}
\DeclareMathOperator*{\Res}{Res}
\newcommand{\abs}[1]{\lvert#1\rvert}
\newcommand{\bbeta}{\boldsymbol\beta}
\newcommand{\dd}{\,\mathrm d}
\newcommand{\e}{\mathrm e}
\newcommand{\ee}{\boldsymbol e}
\newcommand{\fall}[2]{#1^{\underline{#2}}}
\newcommand{\mm}{\boldsymbol m}
\newcommand{\one}{\boldsymbol 1}
\newcommand{\pFq}[5]{{}_{#1}F_{#2}\!\left(\begin{matrix}#3\\#4\end{matrix};#5\right)}
\title{An Askey-Type Confluence Scheme for Hahn-Like Multiple Orthogonality}
\author{Manuel Ma\~nas}
\date{Source: arXiv:2609.08048v1 (CC BY 4.0)}
\begin{document}
\maketitle

\noindent\emph{Source and licence.} An Askey-Type Confluence Scheme for Hahn-Like Multiple Orthogonality. Version arXiv:2609.08048v1 (CC BY 4.0), distributed by arXiv under the Creative Commons Attribution 4.0 International licence, http://creativecommons.org/licenses/by/4.0/, as recorded in the arXiv licence metadata for this exact version and shown on the abstract page https://arxiv.org/abs/2609.08048v1. Full licence notice: \texttt{licenses/arxiv-2609.08048-CC-BY-4.0.txt}.\par\smallskip

\noindent\textit{Editorial scope.} Counted unit: the article's theorem thm:explicit-unreflected-B (source0.tex lines 2538--2573). The article's complete original proof of it is delivered with it (source0.tex lines 2575--2896, one \texttt{proof} environment of 322 lines). No mathematics is written, repaired or re-derived: the statement and the proof are the article's own lines, in the article's own notation, and no retained line is substituted or re-flowed.  Retained uncounted as the source-local context the proof needs: lines 422--433; lines 1229--1234; lines 1253--1259; lines 1261--1276; lines 1390--1397; lines 1399--1411; lines 1541--1561; lines 1625--1630; lines 2191--2237; lines 2241--2251; lines 2258--2262; lines 2295--2318; lines 2323--2344; lines 2346--2362; lines 2368--2373; lines 2378--2385; lines 2405--2414; lines 2435--2441; lines 2444--2451; lines 2469--2476; lines 2487--2492; lines 2506--2515; lines 2523--2530; lines 2532--2536; lines 3194--3200; lines 3212--3265; lines 3333--3379; lines 4245--4250; lines 4260--4288; lines 4332--4374; lines 7187--7191; lines 7201--7281.  Nothing was omitted from any retained span, so the package carries no omission record.  No retained line contains a citation, so the wrapper carries no bibliography and no bibliography entry.  No retained line contains a figure environment, an image-inclusion command or drawing code, so no image is delivered and the package declares no asset dependency.  Editorial formatting only, in addition: an AMS \texttt{amsart} preamble that restates the article's own declarations and macros used by the retained lines, namely the environments \texttt{definition}, \texttt{theorem}, \texttt{corollary} on separate counters, together with the standard packages those lines need.  The retained environments print in this document as Definition 1, Theorem 1, Corollary 1 in order of appearance, whereas the article prints the same items with the numbering of its own full document.  The complete native source of the article as distributed in the arXiv e-print for this exact version is bundled unchanged as \texttt{sources/209691/source0.tex}.
\begin{definition}[Normal multi-index]
\label{def:normal-multi-index}
We call \(\mm\) \emph{normal} for
\((\mathcal L_1,\ldots,\mathcal L_q)\) if \(M_{\mm}\) is nonsingular.
This definition is independent of the chosen monic degree-graded basis.
Equivalently, the normalized type-I problem has a unique solution and the
monic type-II problem has a unique solution of degree \(d\), because their
coefficient matrices are \(M_{\mm}\) and \(M_{\mm}^{\mathsf T}\),
respectively. Thus normality is a simultaneous type-I and type-II
property. Multiplication of any row functional by a nonzero constant does
not affect it.
\end{definition}

\begin{equation}
 \mathcal B_{\mm,N}^{\mathrm H}(k)
 =\sum_{\substack{j\in\{1,\ldots,q\}\\m_j>0}}
 B_{j,N}(k)\widetilde v_{j,N}(k).
 \label{eq:B-component-decomposition}
\end{equation}

\begin{equation}
 \mathcal H^{\mathrm H}_{\mm,N}(z)
 \coloneq\sum_{k=0}^N\mathcal B_{\mm,N}^{\mathrm H}(k)z^k
 =-\fall{N}{n}\frac{(\A)_n}{(\bbeta)_{\mm+n}}(1-z)^n
 \,\pFq{q+1}{q}{n-N,\A+n}{\bbeta+\mm+n}{1-z}
 \label{eq:Hahn-complete-B-pgf}
\end{equation}

\begin{theorem}[Normalized Hahn-like signed sequence]
\label{thm:Hahn-complete-B}
Let \(\mm\) be near the diagonal with
\(1\le d\coloneq\abs\mm\le N+1\), and set \(n=d-1\). The
factorial moments of the sequence defined by
\eqref{eq:Hahn-complete-B-pgf} are
\begin{equation}
 \sum_{k=0}^N\fall{k}{r}\mathcal B_{\mm,N}^{\mathrm H}(k)
 =-h_N(r)\frac{(-r)_n}
 {\prod_{h=1}^q(r+\beta_h)_{m_h}},
 \qquad r\in\Nzero.
 \label{eq:Hahn-complete-B-moments}
\end{equation}
In particular, its first \(n\) moments vanish and its moment of order \(n\)
is nonzero.
\end{theorem}

\begin{equation}
 b_{J,K}=\frac{1}{d_{J,K}}
 \left(\pi_{J,K}-
 \sum_{j=1}^q\sum_{\ell=K+1}^{m_j-1}
 b_{j,\ell}\,
 \Res_{r=-\beta_J-K}R_{j,\ell}^{(N)}(r)\right).
 \label{eq:triangular-B-reconstruction}
\end{equation}

\begin{theorem}[Type-I polynomials from residues]
\label{thm:Hahn-B-components}
Let \(\mm\) be near the diagonal, let \(1\le\abs\mm\le N+1\), and assume
that the pole sets
\(\{-\beta_j-K:K\in\{0,\ldots,m_j-1\}\}\), indexed by
\(j\in\{1,\ldots,q\}\) with \(m_j>0\), are pairwise disjoint and that
\(d_{J,K}\ne0\) for every \(J\in\{1,\ldots,q\}\) with \(m_J>0\) and
every \(K\in\{0,\ldots,m_J-1\}\). Then the coefficients
defined by \eqref{eq:triangular-B-reconstruction} satisfy
\eqref{eq:B-component-decomposition} for the signed sequence of
Theorem~\ref{thm:Hahn-complete-B}, and that sequence is the type-I linear
form for the normalized weights.
\end{theorem}

\begin{corollary}[Closed hypergeometric Hahn-like type-I polynomials]
\label{cor:Hahn-B-closed-components}
Under the hypotheses of Theorem~\ref{thm:Hahn-B-components}, the type-I
polynomials for the rescaled weights
\(\widetilde v_{j,N}=v_{j,N}/\beta_j\) are
\begin{equation}
 B_{j,N}(k)=
 \sum_{J=1}^q\sum_{K=0}^{m_J-1}
 \pi_{J,K}\mathcal C_{j,J,K}^{(N)}(k),
 \qquad j\in\{1,\ldots,q\},\quad m_j>0.
 \label{eq:Hahn-B-closed-components}
\end{equation}
For every \(J\in\{1,\ldots,q\}\) with \(m_J>0\) and every
\(K\in\{0,\ldots,m_J-1\}\),
\(\deg \mathcal C_{J,J,K}^{(N)}\le K\). For every such \(J,K\) with
\(K\ge1\), and every \(j\in\{1,\ldots,q\}\) with \(m_j>0\) and
\(j\ne J\), one has \(\deg \mathcal C_{j,J,K}^{(N)}\le K-1\).
Furthermore, \(\mathcal C_{j,J,0}^{(N)}=0\) for every pair of distinct
active indices \(j,J\), and \(\deg B_{j,N}<m_j\) for every
\(j\in\{1,\ldots,q\}\) with \(m_j>0\).
\end{corollary}

\begin{equation}
 (z+c)I_c(z)=\frac1c+
 \sum_{j=1}^q\alpha_j(c)I_{\beta_j}(z)
 +\alpha_0(c)I_{c+1}(z).
 \label{eq:Hahn-Cauchy-contiguous-recurrence}
\end{equation}

The families use different weight normalizations. Fix a near-diagonal
multi-index \(\mm\) with \(\abs\mm\ge1\); in the Kravchuk case assume
also \(\abs\mm\le N+1\). Throughout this subsection,
\(X\in\{\mathrm K,\mathrm{MII},\mathrm{CII}\}\),
\(j\in\{1,\ldots,q\}\), and \(k,r,\ell,u\in\Nzero\), unless a smaller
range is displayed. Here
\(D_j\) denotes a component relative to a rescaled weight in the
Kravchuk-, Meixner-II-, Charlier-II-, and Laguerre-I-like formulas; the
conversion to the original weight is stated in each case. The symbols
\(B_j\) in the Hahn-, Meixner-I-, Charlier-I-, and Jacobi-like formulas
retain the normalizations fixed in the preceding sections. Put
\(n=\sum_{j=1}^q m_j-1\),
\(D_-(r)=\prod_{h=1}^{q-1}(r+\beta_h)_{m_h}\), and
\(\mathcal R_{\mm}(r)=-(-r)_n/D_-(r)\).
Here \(v_j^X\) and \(\mathcal B_{\mm}^X\) denote the weights and signed
sequences specified in Theorems~\ref{thm:K-rows}--\ref{thm:K-forms},
\ref{thm:MII-rows}--\ref{thm:MII-forms}, and
\ref{thm:CII-system}, respectively. For the
Kravchuk-like family, suppress the fixed lattice size by writing
\(v_j^{\mathrm K}\coloneq v_{j,N}^{\mathrm K}\) and
\(\mathcal B_{\mm}^{\mathrm K}\coloneq\mathcal B_{\mm,N}^{\mathrm K}\), and
extend both sequences by zero to \(k>N\).
For each \(X\in\{\mathrm K,\mathrm{MII},\mathrm{CII}\}\), define the
rescaled weights
\begin{equation}
 w_j^X=\frac{v_j^X}{\beta_j}\quad
 (j\in\{1,\ldots,q-1\}),
 \qquad w_q^X=v_q^X.
 \label{eq:unreflected-canonical-rows}
\end{equation}
Their factorial moments are
\begin{equation}
 \sum_{k\ge0}\fall{k}{r}w_j^X(k)=
 \begin{cases}
 h_X(r)/(r+\beta_j),&j\in\{1,\ldots,q-1\},\\
 h_X(r),&j=q,
 \end{cases}
 \label{eq:unreflected-canonical-moments}
\end{equation}
where
\(h_{\mathrm K}(r)=\fall{N}{r}c^r(\A_{<q})_r/(\bbeta_{<q})_r\),
\(h_{\mathrm{MII}}(r)=\tau^r(\A)_r/(\bbeta_{<q})_r\), and
\(h_{\mathrm{CII}}(r)=a^r(\A_{<q})_r/(\bbeta_{<q})_r\).
Whenever a normalized identity below contains \(h_{\mathrm K}(r)^{-1}\),
it is asserted for \(r\in\{0,\ldots,N\}\). This is the entire range needed
later, since \(n\le N\); the corresponding unnormalized moment identities
remain valid for every \(r\in\Nzero\).

\begin{align}
 R_{j,\ell}^X(r)
 &=\sum_{u=0}^{\ell}\binom\ell u\fall{r}{\ell-u}
   \frac{\Theta_u^X(r)}{r+\beta_j+u},&&
   j\in\{1,\ldots,q-1\},
 \label{eq:unreflected-R-finite-row}\\
 R_{q,\ell}^X(r)
 &=\sum_{u=0}^{\ell}\binom\ell u\fall{r}{\ell-u}
   \Theta_u^X(r),
 \label{eq:unreflected-R-base-row}
\end{align}

\begin{equation}
 \frac1{h_X(r)}\sum_{k\ge0}\fall{k}{r}\fall{k}{\ell}w_j^X(k)
 =R_{j,\ell}^X(r).
 \label{eq:unreflected-R-pairing}
\end{equation}

\begin{align}
 \mathcal C_{j,J,K}^{\mathrm K}(k)
 &=\Delta_{j,J,K}^{\mathrm K}
 \pFq{q+1}{q}
 {1,1-k-\xi_{J,K},
  \bbeta_{<q}+(1-\xi_{J,K})\one_{q-1}-\ee_j^-}
 {1-N-\xi_{J,K},
  \A_{<q}+(1-\xi_{J,K})\one_{q-1}}{c^{-1}},
 \label{eq:K-finite-pole-block}\\
 \mathcal C_{j,J,K}^{\mathrm{MII}}(k)
 &=\Delta_{j,J,K}^{\mathrm{MII}}
 \pFq{q+1}{q}
 {1,1-k-\xi_{J,K},
  \bbeta_{<q}+(1-\xi_{J,K})\one_{q-1}-\ee_j^-}
 {\A+(1-\xi_{J,K})\one_q}{-\tau^{-1}},
 \label{eq:MII-finite-pole-block}\\
 \mathcal C_{j,J,K}^{\mathrm{CII}}(k)
 &=\Delta_{j,J,K}^{\mathrm{CII}}
 \pFq{q+1}{q-1}
 {1,1-k-\xi_{J,K},
  \bbeta_{<q}+(1-\xi_{J,K})\one_{q-1}-\ee_j^-}
 {\A_{<q}+(1-\xi_{J,K})\one_{q-1}}{-a^{-1}}.
 \label{eq:CII-finite-pole-block}
\end{align}

\begin{align}
 \Delta_{j,J,K}^{\mathrm K}
 &=\frac{(N+\beta_j)
   \prod_{h=1}^{q-1}(A_h-\beta_j)
   \prod_{\substack{1\le h\le q-1\\h\ne j}}(\beta_h-\xi_{J,K})}
  {(N+\xi_{J,K})
   \prod_{\substack{1\le h\le q-1\\h\ne j}}(\beta_h-\beta_j)
   \prod_{h=1}^{q-1}(A_h-\xi_{J,K})},
 \label{eq:K-finite-pole-prefactor-row}\\
 \Delta_{j,J,K}^{\mathrm{MII}}
 &=\frac{\prod_{h=1}^{q}(A_h-\beta_j)
   \prod_{\substack{1\le h\le q-1\\h\ne j}}(\beta_h-\xi_{J,K})}
  {\prod_{\substack{1\le h\le q-1\\h\ne j}}(\beta_h-\beta_j)
   \prod_{h=1}^{q}(A_h-\xi_{J,K})},
 \label{eq:MII-finite-pole-prefactor-row}\\
 \Delta_{j,J,K}^{\mathrm{CII}}
 &=\frac{\prod_{h=1}^{q-1}(A_h-\beta_j)
   \prod_{\substack{1\le h\le q-1\\h\ne j}}(\beta_h-\xi_{J,K})}
  {\prod_{\substack{1\le h\le q-1\\h\ne j}}(\beta_h-\beta_j)
   \prod_{h=1}^{q-1}(A_h-\xi_{J,K})}.
 \label{eq:CII-finite-pole-prefactor-row}
\end{align}

\begin{align}
 \Delta_{q,J,K}^{\mathrm K}
 &=(1-c^{-1})
   \frac{\prod_{h=1}^{q-1}(\beta_h-\xi_{J,K})}
   {(N+\xi_{J,K})\prod_{h=1}^{q-1}(A_h-\xi_{J,K})},
 \label{eq:K-finite-pole-prefactor-base}\\
 \Delta_{q,J,K}^{\mathrm{MII}}
 &=-\frac{1+\tau}{\tau}
   \frac{\prod_{h=1}^{q-1}(\beta_h-\xi_{J,K})}
   {\prod_{h=1}^{q}(A_h-\xi_{J,K})},
 \label{eq:MII-finite-pole-prefactor-base}\\
 \Delta_{q,J,K}^{\mathrm{CII}}
 &=-\frac1a
   \frac{\prod_{h=1}^{q-1}(\beta_h-\xi_{J,K})}
   {\prod_{h=1}^{q-1}(A_h-\xi_{J,K})}.
 \label{eq:CII-finite-pole-prefactor-base}
\end{align}

\begin{equation}
 \frac1{r+\xi_{J,K}}
 =\sum_{j=1}^q\frac1{h_X(r)}
   \sum_{k\ge0}\fall{k}{r}\mathcal C_{j,J,K}^X(k)w_j^X(k).
 \label{eq:finite-pole-block-pairing}
\end{equation}

\begin{equation}
 \pi_{J,K}^{-}=
 \frac{(-1)^{K+1}(\beta_J+K)_n}
 {K!(m_J-1-K)!
  \prod_{\substack{1\le h\le q-1\\h\ne J}}
  (\beta_h-\beta_J-K)_{m_h}}.
 \label{eq:unreflected-finite-pole-coefficient}
\end{equation}

\begin{equation}
 \mathfrak h_\nu(\boldsymbol\xi;\boldsymbol\eta)
 =\sum_{\substack{\boldsymbol\epsilon\in\{0,1\}^{p},
                   \,\boldsymbol\kappa\in\Nzero^{t}\\
                   \sum_{i=1}^p\epsilon_i+\sum_{i=1}^t\kappa_i=\nu}}
 (-1)^{\sum_{i=1}^t\kappa_i}
 \prod_{i=1}^{p}\xi_i^{\epsilon_i}
 \prod_{i=1}^{t}\eta_i^{\kappa_i}.
 \label{eq:unreflected-Horn-finite-sum}
\end{equation}

\begin{equation}
 p_{K,\ell}^X=
 \sum_{u=0}^{\ell}\binom\ell u\sigma_u^X
 \mathfrak h_{d_{\ell,u}^X-K}
 (\boldsymbol\Xi_{\ell,u}^X;\boldsymbol\eta_u).
 \label{eq:unreflected-infinity-transition}
\end{equation}

\begin{equation}
 p_{\ell,\ell}^{\mathrm K}=(1-c)^\ell,
 \qquad
 p_{\ell,\ell}^{\mathrm{MII}}=(1+\tau)^\ell,
 \qquad
 p_{\ell,\ell}^{\mathrm{CII}}=1.
 \label{eq:unreflected-infinity-pivots}
\end{equation}

\begin{equation}
 z_{J,K,\ell}^X=
 \sum_{u=K+1}^{\ell}\binom\ell u
 \fall{-\beta_J-K}{\ell-u}
 \frac{W_{J,K,u}^X}
 {(-1)^K K!(u-K-1)!Q_{J,K}(u)}
 \label{eq:unreflected-base-residue}
\end{equation}

\begin{equation}
 \sum_{j=1}^q\frac1{h_X(r)}
 \sum_{k\ge0}\fall{k}{r}\mathcal E_{j,\ell}^X(k)w_j^X(k)
 =\sum_{K=0}^{\ell}p_{K,\ell}^Xr^K.
 \label{eq:unreflected-corrected-infinity-pairing}
\end{equation}

\begin{equation}
 \gamma_\ell^X=
 \sum_{s=0}^{m_q-1-\ell}(-1)^s
 \sum_{\ell=\ell_0<\ell_1<\cdots<\ell_s\le m_q-1}
 \frac{\pi_{\infty,\ell_s}}{p_{\ell_s,\ell_s}^X}
 \prod_{t=0}^{s-1}
 \frac{p_{\ell_t,\ell_{t+1}}^X}
      {p_{\ell_t,\ell_t}^X}.
 \label{eq:unreflected-infinity-chain}
\end{equation}

\begin{equation}
 D_j^X(k)=
 \sum_{J=1}^{q-1}\sum_{K=0}^{m_J-1}
 \pi_{J,K}^{-}\mathcal C_{j,J,K}^X(k)
 +\sum_{\ell=0}^{m_q-1}
 \gamma_\ell^X\mathcal E_{j,\ell}^X(k).
 \label{eq:unreflected-explicit-components}
\end{equation}

\begin{equation}
 B_j^X=D_j^X/\beta_j,\quad j\in\{1,\ldots,q-1\},
 \qquad B_q^X=D_q^X.
 \label{eq:unreflected-probability-components}
\end{equation}

\begin{equation}
 A_q^{(T)}=Ta,\qquad
 \beta_q^{(T)}=Tb,\qquad
 0<a<b,\qquad
 c\coloneq\frac ab\in(0,1),\qquad T\to\infty.
 \label{eq:K-scaling}
\end{equation}

\begin{theorem}[Kravchuk-like weights]
\label{thm:K-rows}
Let \(N\in\Nzero\), use the scaling \eqref{eq:K-scaling}, and assume
\(0<A_h<\beta_h\) for every \(h\in\{1,\ldots,q-1\}\). As
\(T\to\infty\), the Hahn-like weights converge
coefficientwise to normalized strictly positive weights
\(v_{j,N}^{\mathrm K}\), for every \(j\in\{1,\ldots,q\}\). The following
formulas hold.
\begin{enumerate}[label=\textup{(\roman*)}]
\item For every \(j\in\{1,\ldots,q\}\), the generating-function identity
below holds for all \(z\in\mathbb C\), the mass formula holds for every
\(k\in\{0,\ldots,N\}\), and the moment formula holds for every
\(r\in\{0,\ldots,N\}\):
\begin{align}
 G_{j,N}^{\mathrm K}(z)
 &\coloneq\sum_{k=0}^Nv_{j,N}^{\mathrm K}(k)z^k
 =\pFq{q}{q-1}
 {-N,\A_{<q}}{\boldsymbol C_j}{c(1-z)},
 \label{eq:K-pgf}\\
 v_{j,N}^{\mathrm K}(k)
 &=\binom Nk c^k
 \frac{(\A_{<q})_k}{(\boldsymbol C_j)_k}
 \pFq{q}{q-1}
 {k-N,\A_{<q}+k}{\boldsymbol C_j+k}{c},
 \label{eq:K-mass}\\
 \sum_{k=0}^N\fall{k}{r}v_{j,N}^{\mathrm K}(k)
 &=\fall{N}{r}c^r
 \frac{(\A_{<q})_r}{(\boldsymbol C_j)_r},
 \qquad r\in\{0,\ldots,N\}.
 \label{eq:K-factorial-moments}
\end{align}
\item With \(Y=y_1\cdots y_{q-1}\), for every
\(j\in\{1,\ldots,q\}\) and \(z\in\mathbb C\), the generating functions
have the representations below; the mass formula holds for every
\(k\in\{0,\ldots,N\}\):
\begin{align}
 G_{j,N}^{\mathrm K}(z)
 &=
 \int_{(0,1)^{q-1}}
 [1-c(1-z)Y]^N
 \prod_{h=1}^{q-1}b_{A_h,C_{h,j}-A_h}(y_h)
 \dd\boldsymbol y,
 \label{eq:K-positive-pgf}\\
 v_{j,N}^{\mathrm K}(k)
 &=
 \binom Nk
 \int_{(0,1)^{q-1}}
 (cY)^k(1-cY)^{N-k}
 \prod_{h=1}^{q-1}b_{A_h,C_{h,j}-A_h}(y_h)
 \dd\boldsymbol y.
 \label{eq:K-positive-mass}
\end{align}
\end{enumerate}
\end{theorem}

\begin{theorem}[Kravchuk-like type-II polynomial and signed sequence]
\label{thm:K-forms}
As \(T\to\infty\) under the scaling \eqref{eq:K-scaling}:
\begin{enumerate}[label=\textup{(\roman*)}]
\item The Hahn-like type-II polynomials converge coefficientwise to
\(A_{\mm,N}^{\mathrm K}\). This polynomial has degree \(d\), satisfies
\(A_{\mm,N}^{\mathrm K}(0)=1\), and \(\widehat
A_{\mm,N}^{\mathrm K}\) is monic. For every
\(j\in\{1,\ldots,q\}\) with \(m_j>0\),
\begin{equation}
 \sum_{k=0}^N\fall{k}{r}
 A_{\mm,N}^{\mathrm K}(k)v_{j,N}^{\mathrm K}(k)=0,
 \qquad r\in\{0,\ldots,m_j-1\}.
 \label{eq:K-A-orthogonality}
\end{equation}
\item The normalized Hahn-like generating functions satisfy, coefficientwise
in \(z\),
\begin{equation}
 \lim_{T\to\infty}
 (\beta_q^{(T)})_{m_q}
 \mathcal H_{\mm,N}^{\mathrm H}(z)
 =\mathcal H_{\mm,N}^{\mathrm K}(z).
 \label{eq:K-B-limit}
\end{equation}
Its factorial moments are
\begin{equation}
 \sum_{k=0}^N\fall{k}{r}
 \mathcal B_{\mm,N}^{\mathrm K}(k)
 =
 -h_N^{\mathrm K}(r)
 \frac{(-r)_n}
 {\prod_{h=1}^{q-1}(r+\beta_h)_{m_h}},
 \qquad r\in\{0,\ldots,N\}.
 \label{eq:K-B-moments}
\end{equation}
The moments of every order \(r\in\Nzero\) with \(r<n\) vanish, and
\begin{equation}
 \sum_{k=0}^N\fall{k}{n}
 \mathcal B_{\mm,N}^{\mathrm K}(k)
 =
 (-1)^{n+1}n!\,h_N^{\mathrm K}(n)
 \frac{1}{\prod_{h=1}^{q-1}(n+\beta_h)_{m_h}}
 \ne0.
 \label{eq:K-B-leading-moment}
\end{equation}
\end{enumerate}
\end{theorem}

\begin{equation}
 \beta_q^{(N)}=A_q+\frac{N}{\tau},
 \qquad
 \tau=\frac{c}{1-c}>0,
 \label{eq:MII-scaling}
\end{equation}

\begin{theorem}[Meixner-II-like weights]
\label{thm:MII-rows}
As \(N\to\infty\) under \eqref{eq:MII-scaling}, for every
\(j\in\{1,\ldots,q\}\), the Hahn-like weights converge coefficientwise
and in every fixed moment to a normalized positive weight
\(v_j^{\mathrm{MII}}\) on \(\Nzero\) with generating function
\begin{equation}
 G_j^{\mathrm{MII}}(z)
 =\pFq{q}{q-1}{A_1,\ldots,A_q}
 {C_{1,j},\ldots,C_{q-1,j}}{\tau(z-1)}.
 \label{eq:MII-pgf}
\end{equation}
Its factorial moments are, for every \(r\in\Nzero\),
\begin{equation}
 \sum_{k\ge0}\fall{k}{r}v_j^{\mathrm{MII}}(k)
 =\tau^r\frac{\prod_{h=1}^q(A_h)_r}
 {\prod_{h=1}^{q-1}(C_{h,j})_r}.
 \label{eq:MII-factorial-moments}
\end{equation}
The same generating function has the positive integral representation
\begin{equation}
 G_j^{\mathrm{MII}}(z)
 =\int_{(0,1)^{q-1}}
 \frac{\prod_{h=1}^{q-1}b_{A_h,C_{h,j}-A_h}(y_h)
       \dd\boldsymbol y}
 {(1+\tau(1-z)Y)^{A_q}}.
 \label{eq:MII-positive-integral}
\end{equation}
\end{theorem}

\begin{theorem}[Meixner-II-like type-II polynomial and signed sequence]
\label{thm:MII-forms}
Let \(\mm\) be fixed and near the diagonal, and put
\(d=\abs\mm\ge1\) and \(n=d-1\). Under \eqref{eq:MII-scaling}, the
following assertions hold as \(N\to\infty\).
\begin{enumerate}[label=\textnormal{(\roman*)}]
\item The Hahn polynomial converges termwise to
\begin{equation}
 A_{\mm}^{\mathrm{MII}}(k)=
 \pFq{q+1}{q}
 {-d,-k,\beta_1+m_1,\ldots,\beta_{q-1}+m_{q-1}}
 {A_1,\ldots,A_q}{-\tau^{-1}}.
 \label{eq:MII-A}
\end{equation}
This polynomial has exact degree \(d\) and satisfies
\(A_{\mm}^{\mathrm{MII}}(0)=1\).
For every \(j\in\{1,\ldots,q\}\) with \(m_j>0\), it satisfies
\begin{equation}
 \sum_{k\ge0}\fall{k}{r}A_{\mm}^{\mathrm{MII}}(k)
 v_j^{\mathrm{MII}}(k)=0,
 \qquad r\in\{0,\ldots,m_j-1\}.
 \label{eq:MII-A-orthogonality}
\end{equation}

\item The generating function of the scaled signed sequence
\((\beta_q^{(N)})_{m_q}\mathcal B_{\mm,N}^{\mathrm H}\) converges to
\begin{equation}
 \mathcal H_{\mm}^{\mathrm{MII}}(z)
 =-\tau^n\frac{(\A)_n}{(\bbeta_{<q})_{\mm_{<q}+n}}(1-z)^n
 \pFq{q}{q-1}
 {\A+n}{\bbeta_{<q}+\mm_{<q}+n}{\tau(z-1)}.
 \label{eq:MII-B-pgf}
\end{equation}
Its factorial moments are, for every \(r\in\Nzero\),
\begin{equation}
 \sum_{k\ge0}\fall{k}{r}\mathcal B_{\mm}^{\mathrm{MII}}(k)
 =-\tau^r(-r)_n
 \frac{\prod_{h=1}^{q}(A_h)_r}
 {\prod_{h=1}^{q-1}(\beta_h)_{m_h+r}}.
 \label{eq:MII-B-moments}
\end{equation}
\end{enumerate}
\end{theorem}

\begin{equation}
 A_q=R,\qquad \tau_R=\frac aR,
 \qquad c_R=\frac{a}{R+a},\qquad R\to\infty,
 \label{eq:CII-scaling}
\end{equation}

\begin{theorem}[Charlier-II-like system]
\label{thm:CII-system}
As \(R\to\infty\) under \eqref{eq:CII-scaling}, the following assertions hold.
\begin{enumerate}[label=\textnormal{(\roman*)}]
\item For every
\(j\in\{1,\ldots,q\}\), the limiting weight has generating function,
coefficients, and factorial moments
\begin{align}
 G_j^{\mathrm{CII}}(z)
 &=\pFq{q-1}{q-1}{\A_{<q}}{\boldsymbol C_j}{a(z-1)},
 \label{eq:CII-pgf}\\
 v_j^{\mathrm{CII}}(k)
 &=\frac{a^k}{k!}
 \frac{(\A_{<q})_k}{(\boldsymbol C_j)_k}
 \pFq{q-1}{q-1}{\A_{<q}+k}{\boldsymbol C_j+k}{-a},
 \label{eq:CII-mass}\\
 \sum_{k\ge0}\fall{k}{r}v_j^{\mathrm{CII}}(k)
 &=a^r\frac{(\A_{<q})_r}{(\boldsymbol C_j)_r},
 \label{eq:CII-moments}
\end{align}
for \(k,r\in\Nzero\), where
\(\boldsymbol C_j=(\beta_h+\delta_{h,j})_{h=1}^{q-1}\), with no shift
when \(j=q\). With \(Y=y_1\cdots y_{q-1}\), they also satisfy
\begin{align}
 G_j^{\mathrm{CII}}(z)
 &=\int_{(0,1)^{q-1}}\e^{aY(z-1)}
 \prod_{h=1}^{q-1}b_{A_h,C_{h,j}-A_h}(y_h)\dd\boldsymbol y,
 \label{eq:CII-Euler-integral}\\
 v_j^{\mathrm{CII}}(k)
 &=\frac1{k!}\int_{(0,1)^{q-1}}\e^{-aY}(aY)^k
 \prod_{h=1}^{q-1}b_{A_h,C_{h,j}-A_h}(y_h)\dd\boldsymbol y.
 \label{eq:CII-positive-coefficients}
\end{align}
These weights are strictly positive and normalized.

\item For a near-diagonal multi-index \(\mm\) with \(D=\abs\mm\ge1\),
the type-II polynomial is
\begin{equation}
 A_{\mm}^{\mathrm{CII}}(k)
 =\pFq{q+1}{q-1}
 {-D,-k,\bbeta_{<q}+\mm_{<q}}{\A_{<q}}{-a^{-1}}.
 \label{eq:CII-A}
\end{equation}
The polynomial \(A_{\mm}^{\mathrm{CII}}\) has degree \(D\), satisfies
\(A_{\mm}^{\mathrm{CII}}(0)=1\), and, for every
\(j\in\{1,\ldots,q\}\) with \(m_j>0\), obeys
\begin{equation}
 \sum_{k\ge0}\fall{k}{s}A_{\mm}^{\mathrm{CII}}(k)
 v_j^{\mathrm{CII}}(k)=0,
 \qquad s\in\{0,\ldots,m_j-1\}.
 \label{eq:CII-A-orthogonality}
\end{equation}

\item For the same \(\mm\), put \(n=D-1\). The signed sequence has
generating function
\begin{equation}
 \mathcal H_{\mm}^{\mathrm{CII}}(z)
 \coloneq\sum_{k\ge0}\mathcal B_{\mm}^{\mathrm{CII}}(k)z^k
 =-a^n\frac{(\A_{<q})_n}{(\bbeta_{<q})_{\mm_{<q}+n}}(1-z)^n
 \,\pFq{q-1}{q-1}
 {\A_{<q}+n}{\bbeta_{<q}+\mm_{<q}+n}{a(z-1)}.
 \label{eq:CII-B}
\end{equation}
The factorial moments of \(\mathcal B_{\mm}^{\mathrm{CII}}\) satisfy,
for every \(r\in\Nzero\),
\begin{equation}
 \sum_{k\ge0}\fall{k}{r}\mathcal B_{\mm}^{\mathrm{CII}}(k)
 =
 -a^r(-r)_n\frac{(\A_{<q})_r}
 {\prod_{h=1}^{q-1}(\beta_h)_{m_h+r}}.
 \label{eq:CII-B-moments}
\end{equation}
In particular, they vanish for \(r\in\Nzero\) with \(r<n\), whereas
\begin{equation}
 \sum_{k\ge0}\fall{k}{n}\mathcal B_{\mm}^{\mathrm{CII}}(k)
 =(-1)^{n+1}n!a^n
 \frac{(\A_{<q})_n}{\prod_{h=1}^{q-1}(\beta_h)_{m_h+n}}\ne0.
 \label{eq:CII-B-leading-moment}
\end{equation}
\end{enumerate}
\end{theorem}

\begin{theorem}[Kravchuk-, Meixner-II-, and Charlier-II-like type-I components]
\label{thm:explicit-unreflected-B}
Let \(\mm\) be near the diagonal with \(\abs\mm\ge1\), and set
\(n=\abs\mm-1\). In the Kravchuk case assume
\(\sum_{j=1}^q m_j\le N+1\). Assume that the sets
\(\{-\beta_J-K:K\in\Nzero,\ K<m_J\}\), indexed by
\(J\in\{1,\ldots,q-1\}\), are pairwise disjoint and that
the denominators in
\eqref{eq:K-finite-pole-block}--\eqref{eq:CII-finite-pole-prefactor-base},
\eqref{eq:unreflected-finite-pole-coefficient}, and
\eqref{eq:unreflected-base-residue} do not vanish at any of their indicated
indices, apart from the removable diagonal cases with \(K=0\). Then, for
every \(X\in\{\mathrm K,\mathrm{MII},\mathrm{CII}\}\),
\(\deg D_j^X<m_j\) for every
\(j\in\{1,\ldots,q\}\) with \(m_j>0\), and
\begin{equation}
 \mathcal B_{\mm}^X(k)=\sum_{j=1}^qD_j^X(k)w_j^X(k),
 \qquad k\in\Nzero.
 \label{eq:unreflected-component-decomposition}
\end{equation}
For every \(r\in\Nzero\), with the additional restriction \(r\le N\)
when \(X=\mathrm K\),
\begin{equation}
 \sum_{k\ge0}r!\binom{k}{r}\mathcal B_{\mm}^X(k)
 =-h_X(r)\frac{(-r)_n}{D_-(r)}.
 \label{eq:unreflected-complete-moments}
\end{equation}
The moments below order \(n\) vanish, whereas the right-hand side of
\eqref{eq:unreflected-complete-moments} at \(r=n\) is nonzero. Moreover,
the type-I moment matrix is nonsingular. Consequently, \(\mm\) is normal
in the sense of Definition~\ref{def:normal-multi-index}, and
\((D_1^X,\ldots,D_q^X)\) is the unique tuple of type-I polynomials with
the stated degree bounds relative to the rescaled weights
\eqref{eq:unreflected-canonical-rows}. The components relative to the
original weights are given by \eqref{eq:unreflected-probability-components}.
\end{theorem}

\begin{proof}
We first derive the three normalized factorial-moment functions. The falling
factorial product identity
\(\fall{k}{r}\fall{k}{\ell}
=\sum_{u=0}^{\ell}\binom\ell u
\fall{r}{\ell-u}\fall{k}{r+u}\)
contains only \(\ell+1\) terms. The required moment quotients are
\(h_X(r+u)/h_X(r)=\Theta_u^X(r)\).
Substitution in \eqref{eq:unreflected-canonical-moments} gives
\eqref{eq:unreflected-R-finite-row}--
\eqref{eq:unreflected-R-base-row}, and hence
\eqref{eq:unreflected-R-pairing}, for all three families.

We next establish the simple-pole moment identity. For every
\(J\in\{1,\ldots,q-1\}\) and \(K\in\Nzero\) with \(K<m_J\), the
Cauchy-transform recurrence for arbitrary Hahn parameters in the proof of
Corollary~\ref{cor:Hahn-B-closed-components} gives
\begin{equation}
 \sum_{j=1}^q\frac1{h_N(r)}
 \sum_{k=0}^N\fall{k}{r}
 \mathcal C_{j,J,K}^{(N)}(k)
 \widetilde v_{j,N}(k)
 =\frac1{r+\xi_{J,K}}.
 \label{eq:unreflected-Hahn-pairing}
\end{equation}
Indeed, \(\mathcal C_{\cdot,J,K}^{(N)}\) is the coefficient vector of
\(I_{\xi_{J,K}}\) in the iterated relation
\eqref{eq:Hahn-Cauchy-contiguous-recurrence}; the polynomial remainder
has no residue at a lattice point, and the factorial moments of
\(I_{\xi_{J,K}}\) are
\(h_N(r)/(r+\xi_{J,K})\).

For the partial Hahn-to-Kravchuk limit
\eqref{eq:K-scaling}, with \(J\in\{1,\ldots,q-1\}\) and
\(K\in\Nzero\) satisfying \(K<m_J\), every summand of the terminating
hypergeometric polynomial has the factor
\(\frac{(\beta_q^{(T)}+1-\xi_{J,K}-\delta_{j,q})_s}
{(A_q^{(T)}+1-\xi_{J,K})_s}
\xrightarrow[T\to\infty]{}c^{-s}\).
The prefactor has the finite limit
\[
 \left.
 (\beta_q^{(T)})^{-\delta_{j,q}}
 \mathcal D_{j,J,K}^{(N)}
 \right|_{A_q=A_q^{(T)},\,\beta_q=\beta_q^{(T)}}
 \xrightarrow[T\to\infty]{}\Delta_{j,J,K}^{\mathrm K};
\]
for \(j\in\{1,\ldots,q-1\}\) all factors containing \(A_q^{(T)}\) or
\(\beta_q^{(T)}\) cancel in ratios, whereas for \(j=q\) the remaining
factor tends to
\(\lim_{T\to\infty}(N+\beta_q^{(T)})(A_q^{(T)}-\beta_q^{(T)})/
[\beta_q^{(T)}A_q^{(T)}]=1-c^{-1}\).
The factors with indices below \(q\) are precisely those in
\eqref{eq:K-finite-pole-prefactor-row} and
\eqref{eq:K-finite-pole-prefactor-base}.  Hence
\[
 \left.
 (\beta_q^{(T)})^{-\delta_{j,q}}
 \mathcal C_{j,J,K}^{(N)}(k)
 \right|_{A_q=A_q^{(T)},\,\beta_q=\beta_q^{(T)}}
 \xrightarrow[T\to\infty]{}\mathcal C_{j,J,K}^{\mathrm K}(k).
\]

For the Hahn-like-to-Meixner-II-like limit \eqref{eq:MII-scaling}, the factor that
contains the two diverging parameters is instead
\(\frac{(\beta_q^{(N)}+1-\xi_{J,K}-\delta_{j,q})_s}
{(1-N-\xi_{J,K})_s}
\xrightarrow[N\to\infty]{}(-\tau^{-1})^s\).
Moreover,
\[
 \left.
 (\beta_q^{(N)})^{-\delta_{j,q}}
 \mathcal D_{j,J,K}^{(N)}
 \right|_{\beta_q=\beta_q^{(N)}}
 \xrightarrow[N\to\infty]{}\Delta_{j,J,K}^{\mathrm{MII}}.
\]
For \(j=q\), the nonconstant part of this limit is
\[
 \lim_{N\to\infty}
 \frac1{\beta_q^{(N)}}
 \frac{(N+\beta_q^{(N)})
       \prod_{g=1}^{q}(A_g-\beta_q^{(N)})}
      {N\prod_{h=1}^{q-1}(\beta_h-\beta_q^{(N)})}
 =-\frac{1+\tau}{\tau};
\]
the factors independent of \(N\) yield the product quotient in
\eqref{eq:MII-finite-pole-prefactor-base}.  Consequently,
\[
 \left.
 (\beta_q^{(N)})^{-\delta_{j,q}}
 \mathcal C_{j,J,K}^{(N)}(k)
 \right|_{\beta_q=\beta_q^{(N)}}
 \xrightarrow[N\to\infty]{}\mathcal C_{j,J,K}^{\mathrm{MII}}(k).
\]

Finally, in the Meixner-II-like-to-Charlier-II-like scaling
\(A_q=R\), \(\tau=a/R\), the term of order \(s\) contains
\(\frac{(-\tau^{-1})^s}{(A_q+1-\xi_{J,K})_s}
=\frac{(-R/a)^s}{(R+1-\xi_{J,K})_s}
\xrightarrow[R\to\infty]{}(-a^{-1})^s\).
The prefactors indexed by \(j\in\{1,\ldots,q-1\}\) lose the quotient
\((R-\beta_j)/(R-\xi_{J,K})\), whose limit is one. For \(j=q\),
\(-\frac{1+\tau}{\tau(A_q-\xi_{J,K})}
=-\frac{R+a}{a(R-\xi_{J,K})}
\xrightarrow[R\to\infty]{}-\frac1a\).
These are \eqref{eq:CII-finite-pole-prefactor-row} and
\eqref{eq:CII-finite-pole-prefactor-base}, so
\(\mathcal C_{j,J,K}^{\mathrm{MII}}\) converges coefficientwise to
\(\mathcal C_{j,J,K}^{\mathrm{CII}}\).

All the series just considered terminate at \(s=K\) when \(j=J\) and
at \(s=K-1\) when \(j\ne J\).  The limits can therefore be taken in
each finite sum.  For each fixed \(r\), only finitely many factorial
moments enter \eqref{eq:unreflected-Hahn-pairing}. The first two limits
in \eqref{eq:unreflected-Hahn-pairing} and then the
Meixner-II-like-to-Charlier-II-like limit prove
\eqref{eq:finite-pole-block-pairing} for \(X=\mathrm K\),
\(X=\mathrm{MII}\), and \(X=\mathrm{CII}\), respectively. In the
removable case \(K=0\), the polynomial equals one in component \(J\) and
zero in every other component.

We now compute the coefficients required by the signed sequence. Along the
Kravchuk limit,
\[
 \frac{(\beta_q^{(T)})_{m_q}}{(r+\beta_q^{(T)})_{m_q}}
 \left(-\frac{(-r)_n}{D_-(r)}\right)
 \xrightarrow[T\to\infty]{}\mathcal R_{\mm}(r).
\]
Along the Hahn-like-to-Meixner-II-like limit,
\[
 \frac{(\beta_q^{(N)})_{m_q}}{(r+\beta_q^{(N)})_{m_q}}
 \left(-\frac{(-r)_n}{D_-(r)}\right)
 \xrightarrow[N\to\infty]{}\mathcal R_{\mm}(r).
\]
In the subsequent Charlier-II-like limit the
rational factor \(\mathcal R_{\mm}\) contains neither \(A_q\) nor \(\tau\), so
it is unchanged. At \(r=-\beta_J-K\),
\(\prod_{\substack{s\in\{0,\ldots,m_J-1\}\\s\ne K}}(s-K)
=(-1)^K K!(m_J-1-K)!\).
Evaluation of the remaining factors gives
\[
 \Res_{r=-\beta_J-K}\mathcal R_{\mm}(r)
 =\frac{(-1)^{K+1}(\beta_J+K)_n}
 {K!(m_J-1-K)!
  \prod_{\substack{1\le h\le q-1\\h\ne J}}
  (\beta_h-\beta_J-K)_{m_h}}
 =\pi_{J,K}^{-}.
\]
After subtracting these simple fractions, the remainder is a polynomial
of degree at most \(m_q-1\). To determine its coefficients, write
\(z=r^{-1}\) and use
\(-(-r)_n=(-1)^{n+1}r^n\prod_{s=0}^{n-1}(1-sz)\) and
\(D_-(r)=r^{\sum_{h=1}^{q-1}m_h}
\prod_{h=1}^{q-1}\prod_{s=0}^{m_h-1}(1+(\beta_h+s)z)\).
the coefficient of \(r^K\) in the polynomial part is
\[
 (-1)^{n+1}
 \mathfrak h_{m_q-1-K}
 \left(\mathsf f_n;
 ([\beta_1]_{m_1},\ldots,[\beta_{q-1}]_{m_{q-1}})\right)
 =\pi_{\infty,K}.
\]
We have therefore obtained the labeled identity
\begin{equation}
 \mathcal R_{\mm}(r)=
 \sum_{J=1}^{q-1}\sum_{K=0}^{m_J-1}
 \frac{\pi_{J,K}^{-}}{r+\beta_J+K}
 +\sum_{K=0}^{m_q-1}\pi_{\infty,K}r^K.
 \label{eq:unreflected-target-partial-fractions}
\end{equation}
Here every \(\mathfrak h_\nu\) is given by the finite formula
\eqref{eq:unreflected-Horn-finite-sum}.

It remains to construct component vectors for the polynomial part.  For
\(u>K\), the only singular factor in the \(u\)-th summand of
\(R_{q,\ell}^X\) at \(r=-\beta_J-K\) is \((r+\beta_J)_u^{-1}\), whose
residue is \(1/[(-1)^K K!(u-K-1)!]\).
Thus the residues for the three families are
\begin{align*}
 z_{J,K,\ell}^{\mathrm K}
 &=\sum_{u=K+1}^{\ell}\binom\ell u
 \fall{-\beta_J-K}{\ell-u}
 \frac{c^u\fall{N+\beta_J+K}{u}
       \prod_{g=1}^{q-1}(A_g-\beta_J-K)_u}
 {(-1)^K K!(u-K-1)!
  \prod_{\substack{1\le h\le q-1\\h\ne J}}
  (\beta_h-\beta_J-K)_u},\\
 z_{J,K,\ell}^{\mathrm{MII}}
 &=\sum_{u=K+1}^{\ell}\binom\ell u
 \fall{-\beta_J-K}{\ell-u}
 \frac{\tau^u\prod_{g=1}^{q}(A_g-\beta_J-K)_u}
 {(-1)^K K!(u-K-1)!
  \prod_{\substack{1\le h\le q-1\\h\ne J}}
  (\beta_h-\beta_J-K)_u},\\
 z_{J,K,\ell}^{\mathrm{CII}}
 &=\sum_{u=K+1}^{\ell}\binom\ell u
 \fall{-\beta_J-K}{\ell-u}
 \frac{a^u\prod_{g=1}^{q-1}(A_g-\beta_J-K)_u}
 {(-1)^K K!(u-K-1)!
  \prod_{\substack{1\le h\le q-1\\h\ne J}}
  (\beta_h-\beta_J-K)_u}.
\end{align*}
They are the three specializations of
\eqref{eq:unreflected-base-residue}.

For the expansion at infinity, the \(u\)-th summand of
\(R_{q,\ell}^X\) can be written as
\[
 \binom\ell u\sigma_u^X
 r^{d_{\ell,u}^X}
 \frac{\prod_{\xi\in\boldsymbol\Xi_{\ell,u}^X}(1+\xi/r)}
      {\prod_{\eta\in\boldsymbol\eta_u}(1+\eta/r)},
\]
where \(d_{\ell,u}^X\) was defined above. Its coefficient of \(r^K\) is
\(\binom\ell u\sigma_u^X
\mathfrak h_{d_{\ell,u}^X-K}
(\boldsymbol\Xi_{\ell,u}^X;\boldsymbol\eta_u)\).
Summing over \(u\) yields
\eqref{eq:unreflected-infinity-transition}.  At degree \(\ell\), all
values of \(u\) contribute in the Kravchuk-like and Meixner-II-like cases, whereas
only \(u=0\) contributes in the Charlier-II-like case. Hence
\(\sum_{u=0}^{\ell}\binom\ell u(-c)^u=(1-c)^\ell\),
\(\sum_{u=0}^{\ell}\binom\ell u\tau^u=(1+\tau)^\ell\), and
\(p_{\ell,\ell}^{\mathrm{CII}}=1\),
which proves all three assertions in
\eqref{eq:unreflected-infinity-pivots}.

Near-diagonality gives \(\ell\le m_q-1\le m_J\).  Therefore the poles
of \(R_{q,\ell}^X\) are exactly among
\(-\beta_J-K\), with
\(K\in\Nzero\) and \(K\le\min(m_J-1,\ell-1)\). Subtracting
\(z_{J,K,\ell}^X/(r+\xi_{J,K})\) for every such pole leaves the
polynomial part just computed.  Combining this observation with
\eqref{eq:finite-pole-block-pairing} proves
\eqref{eq:unreflected-corrected-infinity-pairing}.

For completeness, set
\(P_\ell^X(r)=\sum_{K=0}^{\ell}p_{K,\ell}^Xr^K\).  Since the diagonal
coefficient in each family is nonzero, the coefficients of an expansion
\(\sum_\ell\gamma_\ell^XP_\ell^X\) must obey
\begin{equation}
 \gamma_\ell^X
 =\frac{\pi_{\infty,\ell}}{p_{\ell,\ell}^X}
 -\frac1{p_{\ell,\ell}^X}
  \sum_{s=\ell+1}^{m_q-1}p_{\ell,s}^X\gamma_s^X.
 \label{eq:unreflected-infinity-recurrence}
\end{equation}
Every increasing index sequence of positive length in
\eqref{eq:unreflected-infinity-chain} has a unique second index \(s\).
Its remaining indices form a sequence occurring in \(\gamma_s^X\).
Grouping the sum by \(s\) turns
\eqref{eq:unreflected-infinity-chain} into
\eqref{eq:unreflected-infinity-recurrence}. Descending from
\(\ell=m_q-1\) proves
\(\sum_{s=\ell}^{m_q-1}p_{\ell,s}^X\gamma_s^X
=\pi_{\infty,\ell}\) for \(0\le\ell<m_q\).
Thus the second term in
\eqref{eq:unreflected-explicit-components} contributes
\(\sum_K\pi_{\infty,K}r^K\) to the normalized moment function, while the
first term contributes the simple fractions in
\eqref{eq:unreflected-target-partial-fractions}. Their sum gives
\(\mathcal R_{\mm}\), proving
\eqref{eq:unreflected-complete-moments}.

This moment identity also proves equality with the signed sequences
specified in the corresponding family sections.  In the Kravchuk case both sides are supported on
\(\{0,\ldots,N\}\), so their factorial moments of orders
\(0,\ldots,N\) determine every coefficient.  In the Meixner-II-like and
Charlier-II-like cases, multiplication of a weight by a polynomial
corresponds in its generating function to applying a polynomial in
\(z\partial_z\). The weight generating functions and the signed generating
functions are analytic in a neighborhood of \(z=1\).  Equation
\eqref{eq:unreflected-complete-moments} says that their Taylor
coefficients at \(z=1\) agree; the identity theorem gives equality of the
generating functions and hence of their coefficients.

The degree restrictions follow directly from the summands. A term
\(\mathcal C_{j,J,K}^X\) has degree at most \(K\) when \(j=J\) and degree
at most \(K-1\) otherwise. In \(\mathcal E_{j,\ell}^X\), the \(q\)-th component
contains the monomial of degree \(\ell<m_q\), and every correction has
\(K<\ell\). If a correction belongs to a component other than \(J\),
near-diagonality gives
\(K-1\le m_J-2\le m_j-1\).  Hence
\(\deg D_j^X<m_j\) for every \(j\in\{1,\ldots,q\}\). Since
\((-r)_n=0\) for every \(r\in\Nzero\) with \(r<n\), the first \(n\)
factorial moments
vanish. At \(r=n\),
\(-h_X(n)\frac{(-n)_n}{D_-(n)}
=(-1)^{n+1}\frac{n!h_X(n)}{D_-(n)}\ne0\), which is the asserted
normalization.

We finish with uniqueness. The space of tuples
\((P_1,\ldots,P_q)\) with \(\deg P_j<m_j\) has dimension
\(\sum_{j=1}^q m_j=n+1\). The component vectors
\[
 \{\mathcal C_{\cdot,J,K}^X:J\in\{1,\ldots,q-1\},\
 K\in\Nzero,\ K<m_J\}
 \ \cup\
 \{\mathcal E_{\cdot,\ell}^X:\ell\in\Nzero,\ \ell<m_q\}
\]
also number \(n+1\). Their normalized moment functions are the distinct simple
fractions \((r+\xi_{J,K})^{-1}\) and the polynomials
\(P_\ell^X\), whose leading coefficients are nonzero.  These rational
functions are linearly independent, so the displayed component vectors
form a basis of the component space.

Suppose a component vector in that space has factorial moments zero for
every \(r\in\{0,\ldots,n\}\), and let \(G(r)\) be its factorial-moment
function divided by \(h_X(r)\).
The formulas for \(R_{j,\ell}^X\) show that \(D_-(r)G(r)\) is a
polynomial. Its degree is at most \(n\): for
\(j\in\{1,\ldots,q-1\}\),
\(R_{j,\ell}^X(r)=\mathrm O(r^{\ell-1})\), whereas
\(R_{q,\ell}^X(r)=\mathrm O(r^\ell)\) as \(r\to\infty\), and
near-diagonality gives the required bounds on \(\ell\).  Since
\(D_-(r)h_X(r)\ne0\) for every \(r\in\{0,\ldots,n\}\), this polynomial has \(n+1\)
distinct zeros and is identically zero. Expansion in the preceding basis
forces every component coefficient to vanish. Thus the kernel of the
square type-I moment map is trivial, so its moment matrix is nonsingular.
By Definition~\ref{def:normal-multi-index}, \(\mm\) is normal and both the
normalized type-I and monic type-II problems are unique.
\end{proof}

\end{document}
